% %%%%%%%%%%%%%%%%%%%%%%%%%%%%%%%%%%%%%%%%%%%%%%%%%%%%%%%%%%%%%%%%%%%%%%%%%%%%%%%%%
%
%	cap#.tex 
%  -----------
%
%	Neste arquivo você escreve o capítulo. Deve sempre iniciar com o ambiente
%	\chapter{Nome do Capitulo}
%
% %%%%%%%%%%%%%%%%%%%%%%%%%%%%%%%%%%%%%%%%%%%%%%%%%%%%%%%%%%%%%%%%%%%%%%%%%%%%%%%%%
%
\chapter{Análises considerando o cenário de túneis gêmeos com galeria transversal}

O presente capítulo introduz o estudo de túneis gêmeos com a construção de uma galeria transversal, e relaciona os mesmos com  túneis gêmeos sem a galeria, apresentados anteriormente. Dessa maneira, é possível analisar o impacto que a construção da galeria transversal causa na distribuição das tensões e na convergência da seção. 

A convergência da seção do túnel é apresentada por $U=-u(r=R_t)/R_t$ e a pressão no revestimento através da expressão $p=-\sigma_{rr}/c$, assim como nos capítulos anteriores. Os resultados são apresentados em termos de final da construção e em longo prazo, ao longo do comprimento do túnel $z/R_t$. As convergências no equilíbrio $U_{eq}$ e as pressões no equilíbrio $p_{eq}$ são dadas na coordenada $z=-140R_t$, ou seja, em uma cota afastada da influência da frente de escavação. As deformações $\rho(\theta)$ e a distribuição das pressões do revestimento $\chi(\theta)$ também serão apresentados ao longo de três seções transversais distintas, para analisar o efeito da galeria.

Serão analisados os mesmos pontos na periferia dos túneis apresentados no capítulo anterior (\autoref{desenho_pontos}), assim como será utilizada a mesma resposta típica no ponto $C(x=-d_1/2,y=R_t,z=-140R_t)$, relativo à coroa do túnel, na exposição dos resultados referentes às pressões no revestimento e convergências ao longo do comprimento dos túneis. Os demais resultados estão apresentados no Anexo B. 

\section{Influência da galeria transversal nos perfis de convergência}

A \autoref{CG_SG_UC_CRE} apresenta os perfis de convergência ao longo do comprimento dos túneis gêmeos com e sem galeria utilizando revestimento elástico. As análises referentes aos túneis gêmeos com galeria utilizando $\phi=10^\circ$ apresentaram problemas de convergência, e por esse motivo não constam nos resultados dessa seção.

Analisando, de modo geral, ambos os estágios final da construção e longo prazo, pode-se observar valores de convergência menores de equilíbrio no cenário sem galeria em comparação aos túneis com galeria. 
(NÃO ENTENDI O MOTIVO)

Observa-se também o efeito do parâmetro friccional do maciço que, quanto maior, tende a diminuir os valores de convergência da seção. Tanto para os túneis com galeria quanto para os túneis sem galeria, nota-se essa redução da convergência ao utilizar $\phi=5^\circ$.

Ao final da construção, os perfis de convergência mais críticos foram observados nos túneis sem galeria, sendo $\phi=0^\circ$ o mais elevado, com $U_{eq} = 1.08\%$. A mesma análise considerando a construção da galeria transversal apresenta uma redução de aproximadamente 6,9\% em relação a esse valor de referência.

As curvas de convergência obtidas a longo prazo possuem uma proximidade considerável das curvas ao final da construção. Esse comportamento é esperado, visto que o revestimento elástico é caracterizado por impedir o desenvolvimento da convergência a longo prazo.

\begin{figure}[H]
	\begin{center}
		\includegraphics[scale=0.8]{CG_SG_UC_CRE_1.pdf}
	\end{center}
	\caption{\label{CG_SG_UC_CRE}Perfis de convergência - túneis gêmeos sem galeria (SG) e com galeria (CG) com revestimento elástico (CRE)}
\end{figure}

Em um cenário onde os túneis gêmeos são analisados com revestimento em concreto viscoelástico, a \autoref{CG_SG_UC_CRVE} demonstra as curvas de convergência para os casos sem galeria e com galeria. Ao analisar as convergências no final da construção do túnel, é possível perceber que os túneis com galeria transversal apresentam resultados de convergência ligeiramente menores em comparação com os túneis sem galeria.

Ao associar o comportamento do revestimento e do maciço, ambos com efeitos viscosos, é possível observar uma evolução das convergências a longo prazo, sendo um aumento de até 46,6\% para $\phi=0^\circ$ dos túneis gêmeos com galeria em comparação com o final da construção.

\begin{figure}[H]
	\begin{center}
		\includegraphics[scale=0.8]{CG_SG_UC_CRVE.pdf}
	\end{center}
	\caption{\label{CG_SG_UC_CRVE}Perfis de convergência - túneis gêmeos sem galeria (SG) e com galeria (CG) com revestimento viscoelástico (CRVE)}
\end{figure}


\section{Influência da galeria transversal nas pressões no revestimento}

A partir dos resultados expostos na \autoref{CG_SG_PC_CRE}, é possível afirmar que as pressões atuantes no revestimento não sofrem mudanças significativas quando comparamos os casos de túneis gêmeos com e sem galeria revestidos com concreto elástico. Para os túneis com galeria transversal, verifica-se uma pequena perturbação das pressões na região de escavação da galeria, nos estágios de final da construção e longo prazo.

Mesmo com um revestimento elástico, devido ao comportamento viscoso do maciço circundante é esperado um leve aumento das pressões atuantes a partir do final da escavação até o estágio final do longo prazo. O maior acréscimo de pressões obtido nas análises resultou em 87,2\%, para o maciço com $\phi=5^\circ$. 

\begin{figure}[H]
	\begin{center}
		\includegraphics[scale=0.8]{CG_SG_PC_CRE.pdf}
	\end{center}
	\caption{\label{CG_SG_PC_CRE}Pressões no revestimento - túneis gêmeos sem galeria (SG) e com galeria (CG) com revestimento elástico (CRE)}
\end{figure}

A \autoref{CG_SG_PC_CRVE} apresenta os resultados em termos de pressão no revestimento viscoelástico de túneis gêmeos com e sem galeria.

\begin{figure}[H]
	\begin{center}
		\includegraphics[scale=0.8]{CG_SG_PC_CRVE.pdf}
	\end{center}
	\caption{\label{CG_SG_PC_CRVE}Pressões no revestimento - túneis gêmeos sem galeria (SG) e com galeria (CG) com revestimento viscoelástico (CRVE)}
\end{figure}

\section{Distribuição anisotrópica das deformações e pressões no revestimento}

Nesta seção, serão apresentados os resultados do estudo da anisotropia da seção transversal em termos de pressões no revestimento e deformações de túneis gêmeos com galeria transversal. Assim, foram adotadas as seções S1, S2 e S3, como demonstra a \autoref{desenho_secoes}, para apresentar os resultados referentes ao estágio final da construção e no longo prazo.

Da mesma forma que no capítulo anterior, a anisotropia das deformações é dada por $\rho(\theta)=U_{eq}(\theta)/U_{eq}(\theta=90^\circ)$, enquanto que a anisotropia das pressões no revestimento é dada por $\chi(\theta)=p_{eq}(\theta)/p_{eq}(\theta=90^\circ)$.

\begin{figure}[H]
	\begin{center}
		\includegraphics[scale=0.8]{desenho_secoes.pdf}
	\end{center}
	\caption{\label{desenho_secoes}Esquema de representação das seções de análise}
\end{figure}

A \autoref{rho_CG_CRE_inst} apresenta os resultados de $\rho(\theta)$ para as seções 1, 2 e 3, considerando túneis gêmeos com revestimento elástico e no estágio final da construção. 

Nessas análises, nota-se que os valores obtidos para os ângulos entre $0^\circ$ e $90^\circ$ são inferiores à referência, ou seja, as deformações na região ao lado do túnel adjacente são menores em comparação com as deformações na coroa. O mesmo comportamento se observa de $\theta=135^\circ$ a $\theta=180^\circ$, região oposta ao túnel adjacente. Isso sugere que as deformações são mais significativas na região superior do túnel, entre os ângulos de $90^\circ$ e $135^\circ$.

A influência do parâmetro friccional é evidente entre os ângulos de $0^\circ$ e $45^\circ$, nas seções S1 e S3. Para o ângulo de atrito $\phi=0^\circ$, observa-se que as deformações nessa região são mais próximas em grandeza das que ocorrem na coroa do túnel, enquanto que um ângulo de atrito maior acarreta em deformações menores.

Os resultados para o longo prazo estão demonstrados na \autoref{rho_CG_CRE_longo}, onde pode-se perceber um comportamento muito similar ao visualizado no final da construção. Há uma predominância de maiores deformações na coroa do túnel, enquanto que os extremos laterais sofrem menos deformações.

\begin{figure}[H]
	\begin{center}
		\includegraphics[scale=0.8]{rho_CG_CRE_inst.pdf}
	\end{center}
	\caption{\label{rho_CG_CRE_inst}Deformada da seção - túneis gêmeos com galeria transversal e revestimento elástico (CRE) no final da construção}
\end{figure}

\begin{figure}[H]
	\begin{center}
		\includegraphics[scale=0.8]{rho_CG_CRE_longo.pdf}
	\end{center}
	\caption{\label{rho_CG_CRE_longo}Deformada da seção - túneis gêmeos com galeria transversal e revestimento elástico (CRE) no longo prazo}
\end{figure}

Os resultados obtidos para a análise com revestimento viscoelástico de túneis gêmeos com galeria transversal estão expostos na \autoref{rho_CG_CRVE_inst}, referente ao final da construção do túnel, e na \autoref{rho_CG_CRVE_longo}, relacionada ao longo prazo.

\begin{figure}[H]
	\begin{center}
		\includegraphics[scale=0.8]{rho_CG_CRVE_inst.pdf}
	\end{center}
	\caption{\label{rho_CG_CRVE_inst}Deformada da seção - túneis gêmeos com galeria transversal e revestimento viscoelástico (CRVE) no final da construção}
\end{figure}

\begin{figure}[H]
	\begin{center}
		\includegraphics[scale=0.8]{rho_CG_CRVE_longo.pdf}
	\end{center}
	\caption{\label{rho_CG_CRVE_longo}Deformada da seção - túneis gêmeos com galeria transversal e revestimento viscoelástico (CRVE) no longo prazo}
\end{figure}

Nas deformações referentes à última escavação, nota-se que em $\theta=0^\circ$, para um ângulo de atrito de $0^\circ$, as deformações são mais elevadas que as observadas nas análises com revestimento elástico, pois o comportamento do revestimento viscoelástico permite que as deformações se desenvolvam lentamente. Ainda no estágio final da construção, pode-se igualmente avaliar a resposta do maciço à variação do ângulo de atrito, conforme já demonstrado anteriormente para o revestimento elástico.

Ao analisar o comportamento em longo prazo, observa-se um incremento significativo das deformações na região superior da seção transversal do túnel, com valores consideravelmente superiores aos registrados em túneis com revestimento de concreto elástico. Este fenômeno decorre das respostas viscosas características tanto do maciço quanto do revestimento.

A \autoref{chi_CG_CRE_inst} apresenta as pressões no revestimento ao longo da seção transversal para túneis gêmeos com galeria utilizando revestimento elástico. É possível verificar que as pressões no revestimento de $\theta=90^\circ$ a $\theta=135^\circ$ tendem a se manter próximas do valor de referência. Por outro lado, em $\theta=0^\circ$ é observado um aumento exorbitante das pressões, chegando a ser pouco mais que o dobro da pressão obtida na coroa do túnel.

Nessa análise das pressões do revestimento, é observado que os resultados nas seções S1 e S3 são relativamente próximos. Essa tendência pode ser visualizada nos demais gráficos a seguir.

\begin{figure}[H]
	\begin{center}
		\includegraphics[scale=0.8]{chi_CG_CRE_inst.pdf}
	\end{center}
	\caption{\label{chi_CG_CRE_inst}Pressões no revestimento - túneis gêmeos com galeria transversal e revestimento elástico (CRE) no final da construção}
\end{figure}

Com os efeitos do longo prazo, a \autoref{chi_CG_CRE_longo} demonstra as pressões atuantes no revestimento elástico. Não se identificam grandes mudanças em relação ao observado no final da escavação, além do aumento dos valores de referência $p_{eq}$. Isso indica que houve um aumento uniforme dessas pressões, ou seja, a relação entre elas nos diferentes ângulos do contorno da seção continua a mesma. 

\begin{figure}[H]
	\begin{center}
		\includegraphics[scale=0.8]{chi_CG_CRE_longo.pdf}
	\end{center}
	\caption{\label{chi_CG_CRE_longo}Pressões no revestimento - túneis gêmeos com galeria transversal e revestimento elástico (CRE) no longo prazo}
\end{figure}

Para o revestimento viscoelástico, observa-se um aumento das pressões em $\theta=0^\circ$, uma tendência já presente no revestimento elástico. Embora a magnitude das pressões atuantes no revestimento tenha diminuído em comparação ao caso elástico, o formato das curvas permaneceu inalterado, evidenciando um comportamento uniforme das pressões nos diferentes cenários analisados. Esse comportamento pode ser visualizado na \autoref{chi_CG_CRVE_inst}, ao final da escavação, e na \autoref{chi_CG_CRVE_longo} que apresenta o longo prazo	.

\begin{figure}[H]
	\begin{center}
		\includegraphics[scale=0.8]{chi_CG_CRVE_inst.pdf}
	\end{center}
	\caption{\label{chi_CG_CRVE_inst}Pressões no revestimento - túneis gêmeos com galeria transversal e revestimento viscoelástico (CRVE) no final da construção}
\end{figure}

\begin{figure}[H]
	\begin{center}
		\includegraphics[scale=0.8]{chi_CG_CRVE_longo.pdf}
	\end{center}
	\caption{\label{chi_CG_CRVE_longo}Pressões no revestimento - túneis gêmeos com galeria transversal e revestimento viscoelástico (CRVE) no longo prazo}
\end{figure}

\section{Relação entre túneis com galeria e sem galeria}

\begin{figure}[H]
	\begin{center}
		\includegraphics[scale=0.8]{CG_razao_SG_U_CRE_0.pdf}
	\end{center}
	\caption{\label{CG_razao_SG_U_CRE_0}Razão entre túneis gêmeos com galeria e sem galeria, com revestimento elástico (CRE) e $\phi = 0^\circ$}
\end{figure}

\begin{figure}[H]
	\begin{center}
		\includegraphics[scale=0.8]{CG_razao_SG_U_CRE_5.pdf}
	\end{center}
	\caption{\label{CG_razao_SG_U_CRE_5}Razão entre túneis gêmeos com galeria e sem galeria, com revestimento elástico (CRE) e $\phi = 5^\circ$}
\end{figure}

\begin{figure}[H]
	\begin{center}
		\includegraphics[scale=0.8]{CG_razao_SG_U_CRVE_0.pdf}
	\end{center}
	\caption{\label{CG_razao_SG_U_CRVE_0}Razão entre túneis gêmeos com galeria e sem galeria, com revestimento viscoelástico (CRVE) e $\phi = 0^\circ$}
\end{figure}

\begin{figure}[H]
	\begin{center}
		\includegraphics[scale=0.8]{CG_razao_SG_U_CRVE_5.pdf}
	\end{center}
	\caption{\label{CG_razao_SG_U_CRVE_5}Razão entre túneis gêmeos com galeria e sem galeria, com revestimento viscoelástico (CRVE) e $\phi = 5^\circ$}
\end{figure}